\section{Symmetries of a cube}

Take a wooden cube, hold it in both hands, turn it, and put it back down in the space it started
in. The cube looks exactly as it did. Whatever you just did to it is a \defn{symmetry} of the cube.
That is the whole idea, and the rest of this section is about taking it seriously enough to count
the symmetries, to name them, and to multiply them together.

\subsection{What counts as a symmetry}

Two restrictions are hidden in the description above, and both matter.

First, the cube was moved \emph{as a solid body}. Distances between its points never changed; it
was not squashed, stretched or bent. A motion of space with that property is called \defn{rigid},
and a symmetry of the cube is a rigid motion that carries the cube onto the same region of space it
occupied before. Because it is rigid, such a motion sends corners to corners, edges to edges and
faces to faces, and it preserves which corners are joined by an edge.

Second, the motion was carried out \emph{with your hands}. You cannot turn a cube into its mirror
image by turning it, however cleverly; you would have to pass it through a mirror. Motions you can
physically perform are the ones that preserve the difference between left-handed and right-handed,
and they are called \defn{orientation-preserving}. Reflections are exactly the motions that fail
this test.

\begin{rem}
Whether to allow reflections is a choice, not a fact, and different books make different choices.
Allowing them doubles every count below: the cube has $24$ rotations and $48$ rigid symmetries
altogether. Throughout this chapter, \emph{symmetry of the cube} means rotation.
\end{rem}

\subsection{Counting the symmetries}

\begin{prop}\label{prop:24}
A cube has exactly $24$ rotational symmetries.
\end{prop}

\begin{proof}
Pick one corner of the cube and call it $A$. A symmetry must send $A$ to a corner, so there are at
most $8$ choices for the image of $A$. Fix that choice. The three edges meeting at $A$ must go to
the three edges meeting at the image corner, in some cyclic order, and a symmetry is completely
determined by where those three edges go. Spinning the cube about the long diagonal through the
image corner cycles those three edges, and there are exactly $3$ positions available. So each of
the $8$ choices admits exactly $3$ symmetries, giving $8\cdot 3=24$, and no symmetry is counted
twice because the image of $A$ together with the spin determines the motion.
\end{proof}

\begin{bigpic}
Knowing that the answer is $24$ is much less useful than it looks. A group is not a number; it is a
set with a multiplication. The interesting question is not how many symmetries there are but
\emph{how they combine}: which turn followed by which turn gives which. That question is what the
rest of the chapter is built to answer, and the answer for the cube turns out to be surprisingly
clean (Theorem~\ref{thm:cubeS4}).
\end{bigpic}

\subsection{The three families of axes}

Every rotation other than ``do nothing'' turns the cube about a line through its centre, and a
useful way to see all $24$ is to sort those lines into three families. In each of the three figures
below, one axis is drawn and the turns about it are shown before and after; a second, smaller
picture shows how many axes of that kind there are.

\subsubsection{Axes through the centres of two opposite faces}

Stand the cube on a table and spin it about the vertical line through the centres of the top and
bottom faces. Quarter turns bring it back to the same region of space, so $90^\circ$, $180^\circ$
and $270^\circ$ are all symmetries; $360^\circ$ is the identity and is counted separately. There
are three such axes, one for each pair of opposite faces.

\sfigs{0.56}{\label{fig:faceaxes}The three face axes, each joining the centres of two opposite
faces, with the six face centres dotted and each axis drawn dashed where it passes through the
inside of the cube. Three axes with three turns about each give $3\cdot3=9$ symmetries.}{\figAxF}

\sfig{\label{fig:facerows}The three turns about one face axis, before and after, with the axis
vertical. Beside each pair is the same turn seen from directly above, looking down the axis: the
square of top corners $D,C,G,H$ with the bottom corners hidden beneath them and the axis as the dot
at the centre.}{\figXbn}

Reading Figure~\ref{fig:facerows} row by row, with corners labelled as in
Figure~\ref{fig:diagonals} below:
\begin{itemize}[nosep,leftmargin=1.3em]
\item \emph{Quarter turn, $90^\circ$.} Each corner of the bottom face moves one place round,
$A\to B\to F\to E\to A$, and the top face does the same, $D\to C\to G\to H\to D$. In the top view
the four corners rotate one quarter of the way round the centre.
\item \emph{Half turn, $180^\circ$.} Every corner jumps to the opposite corner of its face:
$A\leftrightarrow F$, $B\leftrightarrow E$, $C\leftrightarrow H$, $D\leftrightarrow G$. In the top
view each corner swaps with the one diagonally across.
\item \emph{Three-quarter turn, $270^\circ$, which is the same as $-90^\circ$.} The quarter turn
performed the other way: $A\to E\to F\to B\to A$ and $D\to H\to G\to C\to D$. It is the inverse of
the first row, which is why the two are listed as $\pm90^\circ$ rather than as two unrelated moves.
\end{itemize}

\subsubsection{Axes through two opposite corners}

A \defn{long diagonal} joins a corner to the opposite corner, passing through the centre of the
cube. Balance the cube on one corner and spin: a third of a turn, $120^\circ$, brings it back, and
so does $240^\circ$. There are four long diagonals, because the eight corners fall into four
opposite pairs.

\sfigs{0.56}{\label{fig:diagonals}The four long diagonals, each joining two opposite corners through the
centre. Four axes with two turns about each give $4\cdot2=8$ symmetries. The corner labels used
throughout are $A,B,C,D$ in front and $E,F,G,H$ behind.}{\figAxD}

\sfig{\label{fig:diagrows}The turns about the long diagonal $AG$, before and after, with the cube
standing on the corner $A$ so that the axis is vertical. Beside each of the first two pairs is the
view straight down the diagonal from above $G$: a regular hexagon whose six vertices are the six
corners other than $A$ and $G$, alternating between the three next to $G$ and the three next to
$A$. The third row attempts a turn by $60^\circ$ and fails: the turned cube, in colour, does not
sit on the original, in grey.}{\figXdn}

In Figure~\ref{fig:diagrows} the corners $A$ and $G$ never move, since they lie on the axis. The
other six split into two rings of three. The $120^\circ$ turn moves the lower ring one step,
$B\to D\to E\to B$, and the upper ring one step, $C\to H\to F\to C$. The $240^\circ$ turn is the
same motion performed the other way, and the two are inverse to one another.

\begin{why}
\emph{Why $60^\circ$ fails.} The top view is a hexagon, so a sixth of a turn looks plausible. But
the six vertices of that hexagon alternate between two different heights along the axis: three of
them are the corners next to $G$, and three are the corners next to $A$. A sixth of a turn would
carry a corner next to $G$ onto a corner next to $A$, that is, it would move a point to a point at a
different distance along the axis, which no rotation about that axis can do. A third of a turn
carries each ring to itself and is therefore available. In short, the corners of the cube supply
only alternate vertices of the hexagon you think you see.
\end{why}

\sfig{\label{fig:downaxes}Why each kind of axis has the order it does. Looking straight down a face
axis, the outline of the cube is a square, which comes back to itself after a quarter turn. Down a
long diagonal the outline is a regular hexagon, but its six vertices alternate between two heights,
so only a third of a turn is available. Down an edge axis the outline is a rectangle, and only a
half turn works.}{\figwhypi}

\subsubsection{Axes through the midpoints of two opposite edges}

Take an edge and the edge diametrically opposite it, and run a line through both midpoints.
A half turn about that line is a symmetry, and nothing else about it is. The twelve edges make six
opposite pairs, so there are six such axes.

\sfigs{0.46}{\label{fig:edgeaxes}The six edge axes, dashed, with the twelve edge midpoints dotted.
Each axis joins the midpoints of two opposite edges, and the twelve edges make six opposite pairs.
Six axes with one turn about each give $6$ symmetries.}{\figAxE}

\sfig{\label{fig:edgerows}The half turn about the axis through the midpoints of $AB$ and $GH$,
before and after, with the cube standing on the edge $AB$. Every corner swaps with a partner:
$A\leftrightarrow B$, $G\leftrightarrow H$, $C\leftrightarrow E$ and
$D\leftrightarrow F$.}{\figXfn}

Two of these half turns deserve attention later: the axis meets the cube only at two edge
midpoints, so no corner is fixed, and every corner moves. That is what makes this family the one
that produces a single swap of diagonals in Section~\ref{sec:def}.

\subsubsection{The three families together}

\begin{center}
\small
\begin{tabular}{@{}lccc@{}}
\toprule
kind of axis & how many axes & turns about each & symmetries in all\\
\midrule
through two opposite face centres & $3$ & $90^\circ,180^\circ,270^\circ$ & $9$\\
through two opposite corners & $4$ & $120^\circ,240^\circ$ & $8$\\
through two opposite edge midpoints & $6$ & $180^\circ$ & $6$\\
\midrule
and the identity & --- & $0^\circ$ & $1$\\
\midrule
total & & & $24$\\
\bottomrule
\end{tabular}
\end{center}

\noindent The total agrees with Proposition~\ref{prop:24}, which is a genuine check: the two counts
were made in completely different ways.

\sfigs{0.82}{The three kinds of axis in one picture, for reference.}{\figaxes}

\begin{tryit}
Hold a cube, or imagine one. Do the $120^\circ$ turn about a long diagonal and then the $240^\circ$
turn about the same diagonal, and confirm that the cube is back where it started. Then do a quarter
turn about a vertical face axis followed by a quarter turn about a horizontal one, and do the same
two turns in the other order. The results are different. Keep that fact; it is the first evidence
that the multiplication we are about to define is not like multiplication of numbers.
\end{tryit}

\subsection{From turns to permutations}

Here is the observation that turns the cube from a bag of $24$ motions into an algebraic object.

The cube has four long diagonals. Label them $d_1,d_2,d_3,d_4$. A symmetry sends corners to corners
and keeps opposite corners opposite, so it sends each long diagonal to a long diagonal: it produces
a \defn{rearrangement} of the four labels. Rearrangements of four things are much easier to handle
than motions of space, and there are $4!=24$ of them, which is suspicious.

\sfig{The cube with its four long diagonals in colour: $d_1=AG$, $d_2=BH$, $d_3=CE$,
$d_4=DF$. Corners $A,B,C,D$ are in front, $E,F,G,H$ behind.}{\figcube{0.85}}

\begin{thm}\label{thm:cubeS4}
The map that sends a rotation of the cube to the rearrangement it induces on the four long diagonals
is a bijection from the $24$ rotations to the $24$ rearrangements of four objects.
\end{thm}

\begin{proof}
Both sets have $24$ elements, by Proposition~\ref{prop:24} and by $4!=24$, so it is enough to show
that every rearrangement occurs.

Consider first the rearrangements that swap two diagonals and leave the other two alone. Take
$d_1=AG$ and $d_2=BH$. The edge $AB$ joins an endpoint of each, and the opposite edge $GH$ joins
the other two endpoints. The half turn about the axis through the midpoints of $AB$ and $GH$
exchanges $A$ with $B$ and $G$ with $H$, hence exchanges $d_1$ with $d_2$, and one checks that it
fixes $d_3$ and $d_4$. The same argument works for any pair of diagonals, since any two of the four
have a pair of opposite edges joining their endpoints in this way. So all six swaps occur.

Now, any rearrangement of four objects can be achieved by a succession of swaps: move the right
object into the first place, then the right one into the second place, and so on. Performing the
corresponding half turns one after another realises any rearrangement. Hence the map is onto, and
being onto between two sets of the same finite size, it is a bijection.
\end{proof}

\begin{rem}
The proof used the fact that a succession of symmetries is again a symmetry, which is obvious for
physical motions and will be given a name in the next section: \emph{closure}. It also used that
every rearrangement is a succession of swaps, which is proved carefully in
Section~\ref{sec:symmetric}.
\end{rem}

Theorem~\ref{thm:cubeS4} is worth more than the count it repeats. Rearrangements are functions, so
they can be composed mechanically, on paper, with no geometric imagination at all. Once the
dictionary is set up, a question about turning a cube becomes a question about composing two
functions on a four-element set. The table below is that dictionary, family by family; notation for
the right-hand column comes in Section~\ref{sec:symmetric}, and it is worth returning to this table
once that notation is available.

\begin{center}
\small
\begin{tabular}{@{}p{0.40\linewidth}cp{0.30\linewidth}c@{}}
\toprule
rotation of the cube & how many & effect on $d_1,\dots,d_4$ & count\\
\midrule
do nothing & $1$ & leaves all four alone & $1$\\
quarter turn about a face axis, either direction & $6$ & cycles all four & $6$\\
half turn about a face axis & $3$ & two swaps at once & $3$\\
third of a turn about a long diagonal, either direction & $8$ & fixes one, cycles three & $8$\\
half turn about an edge axis & $6$ & swaps two, fixes two & $6$\\
\midrule
total & $24$ & & $24$\\
\bottomrule
\end{tabular}
\end{center}

\begin{why}
\emph{Why a third of a turn about $d_1$ fixes $d_1$ and cycles the other three.} The axis of the
motion is the diagonal $d_1$ itself, so its two endpoints stay put and $d_1$ goes to $d_1$. The
three corners next to $A$ are permuted cyclically by the turn, and the three diagonals starting at
those corners are $d_2,d_3,d_4$, so those three are cycled. This is also why the count in that row
is $8$ and not $4$: each of the four diagonals gives two such turns, one in each direction.
\end{why}

\begin{tryit}
Check the counts in the two right-hand columns of the table line by line. Then take the two half
turns about perpendicular edge axes, perform one after the other, and work out which row of the
table the result belongs to. (You will get a rotation of order three; find its axis.)
\end{tryit}
